% Part: first-order-logic
% Chapter: syntax-and-semantics
% Section: free-vars-sentences

\documentclass[../../../include/open-logic-section]{subfiles}

\begin{document}

\olfileid{fol}{syn}{fvs}

\olsection{Wanneer een \printtoken{P}{variable} vrij is, en wat een
\printtoken{P}{sentence} is}

\begin{defn}[Vrije voorkomens van !!a{variable}]
\ollabel{defn:free-occ}
We bepalen stap voor stap welke voorkomens van !!a{variable} in
!!a{formula} \emph{vrij} zijn. Daarvoor gelden de volgende regels:
\begin{enumerate}
\item \indcase*{!A}{$!A$ geen logische deelopbouw heeft (dus atomair is)}{alle voorkomens van
  !!{variable}s in $\indfrm$ vrij zijn.}

\tagitem{prvNot}{\indcase{!A}{\lnot !B}{de vrije voorkomens van
  !!{variable}s in $\indfrm$ precies de vrije voorkomens uit $!B$ zijn.}}{}

\item \indcase{!A}{(!B \ast !C)}{de vrije voorkomens van
  !!{variable}s in $\indfrm$ bestaan uit de vrije voorkomens in $!B$
  samen met die in~$!C$.}

\tagitem{prvAll}{\indcase{!A}{\lforall[x][!B]}{de vrije voorkomens van
  !!{variable}s in $\indfrm$ precies de vrije voorkomens uit~$!B$ zijn,
  behalve de voorkomens van~$x$.}}{}

\tagitem{prvEx}{\indcase{!A}{\lexists[x][!B]}{de vrije voorkomens van
  !!{variable}s in $\indfrm$ precies de vrije voorkomens uit~$!B$ zijn,
  behalve de voorkomens van~$x$.}}{}
\end{enumerate}
\end{defn}

\begin{defn}[Wanneer een variabelevoorkomen gebonden is]
Een voorkomen van !!a{variable} in een formule~$!A$ is \emph{gebonden}
als het niet vrij is.
\end{defn}

\begin{prob}
Maak nu zelf, stap voor stap, een definitie van de gebonden voorkomens
van variabelen. Dit is een inductieve definitie. Gebruik dezelfde opzet als in
\olref[fol][syn][fvs]{defn:free-occ}.
\end{prob}

\begin{defn}[Het formuledeel waarop een kwantor werkt (bereik)]
\iftag{prvAll}{Stel dat $\lforall[x][!B]$ als formuledeel (subformule) voorkomt in
  een formule~$!A$. Kijk dan naar het bijbehorende voorkomen van~$!B$
  binnen~$!A$. Dat formuledeel heet het \emph{bereik} van het
  bijbehorende voorkomen van~$\lforall[x]$.
  \iftag{prvEx}{Voor $\lexists[x]$ geldt hetzelfde.}{}}{Stel dat
  $\lexists[x][!B]$ als formuledeel (subformule) voorkomt in een formule~$!A$.
  Kijk dan naar het bijbehorende voorkomen van~$!B$ binnen~$!A$. Dat
  formuledeel heet het \emph{bereik} van het bijbehorende voorkomen
  van~$\lexists[x]$.}

Stel dat $!B$ het bereik is van een voorkomen van de kwantor
\iftag{prvAll}{$\lforall[x]$\iftag{prvEx}{ of
    $\lexists[x]$}{}}{$\lexists[x]$} in~$!A$. Dan zijn de voorkomens
van $x$ die binnen~$!B$ nog vrij zijn, gebonden in
\iftag{prvAll}{$\lforall[x][!B]$\iftag{prvEx}{ en
$\lexists[x][!B]$}{}}{$\lexists[x][!B]$}. We zeggen daarom dat het
genoemde voorkomen van de kwantor deze voorkomens
\emph{bindt}.
\end{defn}

\begin{ex}
Kijk naar de volgende formule:
\[
\lexists[\Obj v_0][\underbrace{\Atom{\Obj A^2_0}{\Obj v_0,\Obj v_1}}_{!B}] 
\]
$!B$ is het bereik van $\lexists[\Obj v_0]$.
De kwantor bindt het voorkomen van $\Obj v_0$ in $!B$. Het voorkomen
van $\Obj v_1$ bindt hij niet. $\Obj v_1$ is hier dus een vrije
variabele.


In een ingewikkelder !!{formula}~$!A$ werkt het op dezelfde manier:
\[
\lforall[\Obj v_0][\underbrace{(\Atom{\Obj A^1_0}{\Obj v_0} \lif
    \Atom{\Obj A^2_0}{\Obj v_0, \Obj v_1})}_{!B}] \lif \lexists[\Obj
  v_1][\underbrace{(\Atom{\Obj A^2_1}{\Obj v_0, \Obj v_1} \lor \lforall[\Obj v_0][\overbrace{\lnot \Atom{\Obj A^1_1}{\Obj v_0}}^{!D}])}_{!C}]
\]
$!B$ is het bereik van de eerste $\lforall[\Obj v_0]$. $!C$ is het
bereik van $\lexists[\Obj v_1]$. $!D$ is het bereik van de tweede
$\lforall[\Obj v_0]$. De eerste $\lforall[\Obj v_0]$ bindt de voorkomens van $\Obj v_0$
in~$!B$. $\lexists[\Obj v_1]$ bindt het voorkomen van $\Obj v_1$ in
$!C$. De tweede $\lforall[\Obj v_0]$ bindt het voorkomen van
$\Obj v_0$ in~$!D$. Het eerste voorkomen van $\Obj v_1$ en het vierde voorkomen van
$\Obj v_0$ zijn vrij in~$!A$. Het laatste voorkomen van $\Obj v_0$ is
vrij als je alleen naar $!D$ kijkt. Binnen $!C$ en binnen de hele
formule~$!A$ is datzelfde voorkomen wel gebonden.
\end{ex}

\begin{defn}[Gesloten formule (zin)]
!!^a{formula}~$!A$ is \article{sentence} \emph{!!{sentence}} dan en
slechts dan als er geen vrije voorkomens van !!{variable}s in staan.
\end{defn}

% add examples!
\end{document}
